The long-only minimum variance portfolio is routinely solved by assembling an $n\times n$
sample covariance matrix and dispatching it to a general-purpose quadratic program. Both
steps are unnecessary: the covariance matrix is not a primitive of the problem.
It is a computational detour.

We propose a Krylov-based solver that works directly on the $T\times n$ returns matrix $X$
without ever forming $\hat\Sigma = X^\top X/T$. Conjugate gradients applied to the SPD
system $\hat\Sigma\,v = \mathbf{1}$ evaluates each matrix-vector product
as two $O(nT)$ operations with $X$; non-negativity is enforced by a primal-dual
active-set outer loop that restarts CG on successively smaller subproblems.
The $O(\sqrt\kappa)$ iteration count of CG, where $\kappa = \kappa(\hat\Sigma)$ is the
spectral condition number, governs runtime and contrasts sharply with the $O(\kappa)$
rate of first-order projected-gradient methods that do not exploit Krylov structure.

Ledoit-Wolf shrinkage~\cite{ledoit2004}, already standard practice for statistical
reasons, turns out to be identical to spectral preconditioning of the CG system: the
same ridge term that shrinks sample eigenvalues toward their mean also compresses
$\kappa$ and directly reduces CG iterations. On S\&P~500 equity data this brings
the iteration count from 744 (sample covariance) to 82 (Ledoit-Wolf).
Statistical regularization and spectral preconditioning are the same operation applied
to the same matrix---a connection that has not been made explicit in the existing
literature.

We benchmark against interior-point references (Clarabel, OSQP), a direct KKT
factorization, and a simplex-projection proximal gradient method included as
a non-Krylov first-order comparator. Experiments on S\&P~500 equity data and
synthetic problems of growing size confirm order-of-magnitude speedups over
interior-point solvers. For standard portfolio sizes the direct KKT solve already
captures most of the gain; at larger $n$ matrix-free CG becomes decisive, with the
$\sqrt\kappa$ advantage over proximal gradient widening as the problem scales.
